The SAT solution was made using two different modeling approaches. Both are included in the report. 

\vspace{1em}
\textit{Model~1} was implemented using several alternative SAT encodings for the cardinality
constraints introduced above. In particular, we considered the \textbf{Bitwise} encoding,
the \textbf{Heule} encoding, the \textbf{Pairwise} encoding, and the \textbf{Sequential} encoding.
Each encoding provides a different trade-off between the number of variables,
the number of clauses, and the strength of constraint propagation.

\vspace{1em}
\textit{Model~2} extends the formulation of Model~1 by adding further constraints and auxiliary decision variables. In contrast to Model~1, the implementation adopts a single SAT encoding approach (e.g., pairwise-style cardinality constraints) and does not perform a systematic comparison across alternative encodings.  

Moreover, the implementation can optionally enable a \textbf{precomputing} phase described in paragraph \ref{par:precomputing_sat}.





\subsection{Decision variables}

\subsubsection{Model - 1}

The use of multiple encodings within the same model allows for a systematic
comparison of their impact on solver performance, while keeping the underlying
problem formulation fixed.

\paragraph{Booleans}

\begin{itemize}
    \item  $ X_{t,h,p,w}\in\{0,1\} \quad\text{meaning team }t\text{ is scheduled with role }h\text{ in slot }(p,w). $
\end{itemize}

\subsubsection{Model - 2}

\paragraph{Booleans}

\begin{itemize}
    \item  $ M_{t_1,t_2,w} \in \{0,1\} \quad\text{meaning teams $t_1$ and $t_2$ play against each other in week $w$}. $
    \item $HOME_{t,w} \in \{0,1\} \quad\text{meaning team $t$ plays at home in week $w$ (away otherwise)}.$
    \item $P_{t,p,w} \in \{0,1\} \quad\text{meaning team $t$ is assigned to period $p$ in week $w$}.$
\end{itemize}


\paragraph{Remarks.}

Model--2 introduces explicit pair variables \(M_{t_1,t_2,w}\) to capture
which two teams meet in a given week, home assignment variables
\(HOME_{t,w}\) to distinguish roles, and period variables
\(P_{t,p,w}\) to ensure consistency of scheduling.


\subsection{Objective function}

\subsubsection{Model - 1}

We minimize total home/away imbalance across teams:
\[
%\min
\;\; \sum_{t=1}^{T} \left| 
\sum_{p=1}^{P}\sum_{w=1}^{W} \big( X_{t,0,p,w} - X_{t,1,p,w} \big)
\right|.
\]
In the implementation with \texttt{Z3}, the optimization is performed by 
\emph{iterative (binary) search} on a pseudo-Boolean bound \(z\). 
At each step we introduce a constraint of the form 
\(\sum_t \lvert \cdot \rvert \leq z\), encoded using cardinality 
constraints, and solve for feasibility. If satisfiable, the bound 
is tightened; otherwise the process stops at the last feasible value.

\subsubsection{Model - 2 }

As in Model--1, the optimization criterion is to balance the number of
home and away games for each team. Formally:

\[
%\min 
\sum_{t=1}^T \left| \sum_{w=1}^W HOME_{t,w}
      - \Big(W - \sum_{w=1}^W HOME_{t,w}\Big) \right|.
\]

% ======= CONSTRAINT =======
\subsection{Constraints}

\subsubsection{Model - 1 }

\paragraph{Main constraints.}

\begin{itemize}
\item[(C1)] \textbf{Every pair of teams plays at most once:}  
%For each distinct pair of teams $t_1,t_2 \in \{1,\dots,T\}$:
\[
\forall\, t_1,t_2 \;\;\text{with}\;\;t_1\neq t_2\quad
\sum_{p=1}^{P}\sum_{w=1}^{W}
\Big( \bigvee_{h\in\{0,1\}} X_{t_1,h,p,w} \;\land\;
       \bigvee_{h\in\{0,1\}} X_{t_2,h,p,w} \Big)
\;\;\leq 1.
\]

\item[(C2)] \textbf{Each team plays exactly once per week:}  
%For every team $t$ and week $w$:
\[
\forall \,t, w\;
\sum_{h\in\{0,1\}}\sum_{p=1}^{P} X_{t,h,p,w} \;=\; 1.
\]

\item[(C3)] \textbf{Each team appears at most twice in the same period over the tournament:}  
%For every team $t$ and period $p$:
\[
\forall\, t, p\;
\sum_{h\in\{0,1\}}\sum_{w=1}^{W} X_{t,h,p,w} \;\leq\; 2.
\]
\end{itemize}

\paragraph{Implied constraints.}

\begin{itemize}
\item[(C4)] \textbf{Each game slot has exactly one home and one away team:}  
%For every period $p$ and week $w$:
\[
\forall\, p, w\quad
\sum_{t=1}^{T} X_{t,1,p,w} = 1
\quad\text{and}\quad
\sum_{t=1}^{T} X_{t,0,p,w} = 1.
\]
\end{itemize}

\subsubsection{Model - 2 }

\paragraph{Main constraints.}

\begin{itemize}
\item[(C1)] \textbf{Each team plays with every other team only once;}
\[
\forall\, t_1<t_2\qquad
\sum_{w=1}^{W} M_{t_1,t_2,w} \;=\; 1.
\]

\item[(C2)] \textbf{Each team plays exactly once per week}
\[
\forall\, t,\,w\qquad
\sum_{p=1}^{P} P_{t,p,w} \;=\; 1.
\]

\item[(C3)] \textbf{Each team plays at most twice in the same period}
\[
\forall\, t,\,p\qquad
\sum_{w=1}^{W} P_{t,p,w} \;\le\; 2.
\]

\item[(C4)] \textbf{Each slot \((p,w)\) hosts exactly two teams.}
\[
\forall\, p,\,w\qquad
\sum_{t=1}^{T} P_{t,p,w} \;=\; 2.
\]
\end{itemize}

\paragraph{Implied constraints.}

\begin{itemize}
\item[(C5)] \textbf{Home/away consistency when two teams meet in week \(w\).}
\[
\forall\, w,\, t_1<t_2\qquad
M_{t_1,t_2,w} \;\Rightarrow\; \big(HOME_{t_1,w} \oplus HOME_{t_2,w}\big).
\]

\item[(C6)] \textbf{Link between pair variables and period assignments.}
\[
\forall\, w,\, t_1<t_2\qquad
M_{t_1,t_2,w} \;\iff\; \bigvee_{p=1}^{P}\big(P_{t_1,p,w}\land P_{t_2,p,w}\big),
\]
%\[\forall\, w,\,p,\, t_1<t_2\qquad\big(P_{t_1,p,w}\land P_{t_2,p,w}\big)\Rightarrow\; M_{t_1,t_2,w}.\]
%(Equivalently: \(M_{t_1,t_2,w}\) iff the two teams share some period \(p\) in week \(w\).)
\end{itemize}

\paragraph{Precomputing tecnique.}
\label{par:precomputing_sat}

\begin{itemize}
\item[(P1)] \textbf{Precomputed round-robin pruning.}
Let \(\mathcal{R}_w\subseteq\{\{t_1,t_2\}\mid t_1<t_2\}\) be the fixed set of pairs for week \(w\) produced by the round-robin generator. Then:
\[
\forall\, w,\, t_1<t_2\qquad
M_{t_1,t_2,w} =
\begin{cases}
1, & \text{if }\{t_1,t_2\}\in \mathcal{R}_w,\\[4pt]
0, & \text{otherwise.}
\end{cases}
\]
\end{itemize}


% ======= VALIDATION SAT =======

\subsection{Validation}
\label{sub:validation_sat}

\paragraph{Experimental Setup} \mbox{}\\
\\
All experiments were conducted using the \texttt{Z3} solver for SAT and
optimization (\cite{z3}). The implementation is provided in Python, and the solver is
invoked through the \textbf{Z3 API}. 
\\
%The experiments were tested on \emph{Linux} on a laptop equipped with a
%\textit{12th Gen Intel(R) Core(TM) i7-1280P} processor and an
%\textit{NVIDIA RTX-3060 GPU}. No GPU acceleration was employed for the SAT
%solving itself, which ran entirely on the CPU.

Reproducibility instructions, including installation steps and execution
commands, are provided in the accompanying \texttt{README.md} file.



\paragraph{Experimental results} \mbox{}\\
\label{par:Experimental_result_SAT}

The \textbf{following tables reports} the results of the SAT experiments. 
Each row corresponds to a different value of $n$. 
On the tables are reported tthe value of the objective value obtained using different approach and in brackets the computation time for the best solution found. If the instance is solved to optimality is in bold. In addition if the instance is proved to be
unsatisfiable, it is indicate as \texttt{UNSAT} and if no answer is obtained within the
time limit of 300s are indicate it with \texttt{N/A}.
Each modeling approach have is own table (Tab.~\ref{tab:results_sat_m1} and Tab.~\ref{tab:results_sat_m2}).  

The \textbf{precomputation time} required to generate the round-robin pairs is not reported in the tables, as it is negligible: it remains below one second even for 
$n\approx50$, and therefore has no practical impact on the overall computational performance.

% MODEL 1 and MODEL 2 - UPDATED TABLE
\begin{table}[h]
\centering
\scriptsize

\begin{minipage}[t]{0.55\textwidth}
\centering
\setlength{\tabcolsep}{3pt}
\renewcommand{\arraystretch}{1.05}
\begin{tabular}{c|c|c|c|c}
\hline
$n$ & Bitwise & Heule & Pairwise & Sequential \\
\hline
2   & \textbf{2} (0s)  & \textbf{2} (0s)  & \textbf{2} (0s)  & \textbf{2} (0s) \\
4   & \texttt{UNSAT} (0s) & \texttt{UNSAT} (0s) & \texttt{UNSAT} (0s) & \texttt{UNSAT} (0s) \\
6   & \textbf{6} (0s)  & \textbf{6} (0s)  & \textbf{6} (0s)  & \textbf{6} (0s) \\
8   & 14 (1s)          & 24 (1s)          & 18 (1s)          & 24 (0s) \\
10  & 26 (6s)          & 36 (4s)          & 18 (7s)          & 20 (6s) \\
12  & 30 (87s)         & 44 (13s)         & 36 (40s)         & 30 (151s) \\
14  & \texttt{N/A}     & \texttt{N/A}     & \texttt{N/A}     & \texttt{N/A} \\
16  & \texttt{N/A}     & \texttt{N/A}     & \texttt{N/A}     & \texttt{N/A} \\
\hline
\end{tabular}
\caption{Experimental results of Model 1.}
\label{tab:results_sat_m1}
\end{minipage}
\hfill
\begin{minipage}[t]{0.42\textwidth}
\centering
\setlength{\tabcolsep}{3pt}
\renewcommand{\arraystretch}{1.05}
\begin{tabular}{c|c|c}
\hline
$n$ & Standard & Precomputing \\
\hline
2   & \textbf{2} (--)      & \textbf{2} (0s) \\
4   & \texttt{UNSAT} (--)  & \texttt{UNSAT} (0s) \\
6   & \textbf{6} (--)      & \textbf{6} (0s) \\
8   & \textbf{8} (--)      & \textbf{8} (0s) \\
10  & \textbf{10} (--)     & \textbf{10} (0s) \\
12  & \textbf{12} (3s)     & \textbf{12} (1s) \\
14  & \textbf{14} (83s)    & \textbf{14} (6s) \\
16  & 46 (244s)            & \textbf{16} (10s) \\
18  & \texttt{N/A}         & 44 (53s) \\
20  & \texttt{N/A}         & 200 (272s) \\
22  & \texttt{N/A}         & \texttt{N/A} \\
\hline
\end{tabular}
\caption{Experimental results of Model 2.}
\label{tab:results_sat_m2}
\end{minipage}
\end{table}
